
% we include the glossary here (frontmatter is included with \input, so this command is as if it was in main.tex)
%All acronyms must be written in this file.
\newacronym{SDV}{SDV}{Software Defined Vehicles}
\newacronym{FFI}{FFI}{Foreign Function Interfaces}
\newacronym{ECU}{ECU}{Electronic Control Units}
\newacronym{RTS}{RTS}{Real-Time System}
\newacronym{GPOS}{GPOS}{General Purpose Operating System}
\newacronym{RTOS}{RTOS}{Real-Time Operating System}
\newacronym{PGF}{PGF}{Portable Graphics Format}
\newacronym{RAII}{RAII}{Resource Acquisition Is Initialization}
\newacronym{STD}{STD}{Standard C++ library}
\newacronym{TMP}{TMP}{Template Metaprogramming}
\newacronym{CWE}{CWE}{Common Weakness Enumeration}
\newacronym{SCORE}{SCORE}{Safety-Critical Open Research Environment}
\newacronym{ASIL}{ASIL}{Automotive Safety Integrity Levels}
\newacronym{QM}{QM}{Quality Management}
\newacronym{WCET}{WCET}{Worst Case Execution Time}


\newacronym{X}{X}{X}

\frontmatter % Use roman page numbering style (i, ii, iii, iv...) for the pre-content pages

\pagestyle{plain} % Default to the plain heading style until the thesis style is called for the body content

%----------------------------------------------------------------------------------------
%	TITLE PAGE
%----------------------------------------------------------------------------------------

\maketitlepage


%----------------------------------------------------------------------------------------
%	STATEMENT of INTEGRITY
%----------------------------------------------------------------------------------------

\integritystatement


%----------------------------------------------------------------------------------------
%	DEDICATION  (optional)
%----------------------------------------------------------------------------------------
%
%\dedicatory{For/Dedicated to/To my\ldots}
% \begin{dedicatory}

%The dedicatory is optional. 

% \end{dedicatory}


%----------------------------------------------------------------------------------------
%	ABSTRACT PAGE
%----------------------------------------------------------------------------------------

\begin{abstract}

The growing complexity of safety-critical software systems, particularly in the automotive domain, 
has intensified concerns regarding memory safety, concurrency correctness 
and long-term maintainability of large C++ codebases. Although C++ remains the dominant language in embedded and 
real-time systems due to its performance and mature ecosystem, 
its reliance on manual memory management continues to be a major source of vulnerabilities. 
Rust has emerged as a promising alternative, offering strong compile-time guarantees for memory safety 
and data-race freedom without introducing garbage collection. 
However, the adoption of Rust in industrial environments is hindered by the scale of existing C++ systems 
and their deep dependence on template-heavy frameworks such as the Boost C++ Libraries.

This dissertation investigates strategies for incrementally modernizing legacy C++ systems by integrating Rust in a safe 
and performant manner. Focusing on Boost-based components commonly used in automotive software architectures, 
it proposes a systematic methodology for selecting migration candidates, designing bidirectional interoperability interfaces, 
and re-implementing critical functionality in Rust. The work explores \gls{FFI} patterns, 
evaluates existing interoperability tools, and addresses the architectural challenges arising from the mismatch 
between C++ templates and Rust’s ownership model.

A Rust-based reimplementation of selected Boost components is developed and integrated into a hybrid C++/Rust system. 
Performance, memory usage, and interoperability overhead are experimentally evaluated to quantify the benefits 
and costs of this approach. The results demonstrate that incremental migration can significantly improve memory safety 
while preserving determinism and performance, providing practical guidance for the adoption of Rust in high-assurance 
and safety-critical software systems.

\end{abstract}

\begin{abstractotherlanguage}
A crescente complexidade dos sistemas de software crítico e seguro, particularmente no domínio automóvel, 
intensificou as preocupações relativas à segurança da memória, à correção da concorrência 
e à manutenção a longo prazo de grandes bases de código C++. Embora o C++ continue a ser a linguagem dominante 
em sistemas embebidos e de tempo real devido ao seu desempenho e ecossistema maduro, 
a sua dependência da gestão manual da memória continua a ser uma importante fonte de vulnerabilidades. 
O Rust surgiu como uma alternativa promissora, oferecendo fortes garantias em tempo de compilação para a segurança da memória 
e a ausência de conflitos de dados, sem introduzir a recolha de lixo. 
No entanto, a adoção do Rust em ambientes industriais é dificultada pela escala dos sistemas C++ existentes 
e sua profunda dependência de frameworks fortemente baseadas em modelos, como as bibliotecas Boost C++.

Esta dissertação investiga estratégias para modernizar incrementalmente sistemas C++ antigos, integrando Rust de maneira segura 
e com bom desempenho. Com foco em componentes baseados em Boost comumente usados em arquiteturas de software automotivo, 
ela propõe uma metodologia sistemática para selecionar candidatos à migração, projetar interfaces de interoperabilidade 
bidirecionais e reimplementar funcionalidades críticas em Rust. O trabalho explora padrões de \gls{FFI}, 
avalia ferramentas de interoperabilidade existentes e aborda os desafios arquitetônicos decorrentes da incompatibilidade 
entre modelos C++ e o modelo de propriedade do Rust.

Uma reimplementação baseada em Rust de componentes Boost selecionados é desenvolvida e integrada a um sistema híbrido C++/Rust. 
O desempenho, o uso de memória e a sobrecarga de interoperabilidade são avaliados experimentalmente para quantificar os benefícios 
e custos dessa abordagem. Os resultados demonstram que a migração incremental pode melhorar significativamente a segurança da memória ,
preservando o determinismo e o desempenho, fornecendo orientações práticas para a adoção do Rust em sistemas de software de alta garantia 
e segurança crítica.

\end{abstractotherlanguage}


%----------------------------------------------------------------------------------------
%	ACKNOWLEDGEMENTS (optional)
%----------------------------------------------------------------------------------------

% \begin{acknowledgements}

%Acknowledgements (optional) to family, supervisors, colleagues, organizations, funding agencies, …

% \end{acknowledgements}


%----------------------------------------------------------------------------------------
%	LIST OF CONTENTS/FIGURES/TABLES PAGES
%----------------------------------------------------------------------------------------

\tableofcontents % Prints the main table of contents

\listoffigures % Prints the list of figures

\listoftables % Prints the list of tables

%\iflanguage{portuguese}{
%\renewcommand{\listalgorithmname}{Lista de Algor\'itmos}
%}
%\listofalgorithms % Prints the list of algorithms
%\addchaptertocentry{\listalgorithmname}


%\renewcommand{\lstlistlistingname}{List of Source Code}
%\iflanguage{portuguese}{
%\renewcommand{\lstlistlistingname}{Lista de C\'odigo}
%}
%\lstlistoflistings % Prints the list of listings (programming language source code)
%\addchaptertocentry{\lstlistlistingname}


%----------------------------------------------------------------------------------------
%	ABBREVIATIONS
%----------------------------------------------------------------------------------------
\begin{abbreviations}{ll} % Include a list of abbreviations (a table of two columns)
%\textbf{LAH} & \textbf{L}ist \textbf{A}bbreviations \textbf{H}ere\\
%\textbf{WSF} & \textbf{W}hat (it) \textbf{S}tands \textbf{F}or\\
\end{abbreviations}

%----------------------------------------------------------------------------------------
%	SYMBOLS
%----------------------------------------------------------------------------------------

%\begin{symbols}{lll} % Include a list of Symbols (a three column table)

% [Note: Although acronyms and symbols are defined in this section, they should also be defined at least the first time used in the dissertation body.]

%$a$ & distance & \si{\meter} \\
%$P$ & power & \si{\watt} (\si{\joule\per\second}) \\
%Symbol & Name & Unit \\

%\addlinespace % Gap to separate the Roman symbols from the Greek

%$\omega$ & angular frequency & \si{\radian} \\

%\end{symbols}



%----------------------------------------------------------------------------------------
%	ACRONYMS
%----------------------------------------------------------------------------------------

\newcommand{\listacronymname}{List of Acronyms}
\iflanguage{portuguese}{
\renewcommand{\listacronymname}{Lista de Acr\'onimos}
}

%Use GLS
\glsresetall
\printglossary[title=\listacronymname,type=\acronymtype,style=long]

%----------------------------------------------------------------------------------------
%	DONE
%----------------------------------------------------------------------------------------

\mainmatter % Begin numeric (1,2,3...) page numbering
\pagestyle{thesis} % Return the page headers back to the "thesis" style
